\documentclass[12pt, a4paper]{ctexart}

% 设置页面边距
\usepackage{geometry}
\geometry{left=2.5cm, right=2.5cm, top=2.5cm, bottom=2.5cm}

% 用于自定义列表格式
\usepackage{enumitem}

% 用于绘制专业的三线表
\usepackage{booktabs}

% 取消段落首行缩进，更适合试卷排版
\setlength{\parindent}{0pt}

\begin{document}

% 标题部分
\begin{center}
    \LARGE \textbf{《AI2803-计算机体系结构》作业1}
    \\
    \large \textbf{姓名：韩岳成 \quad 学号：524531910029}
\end{center}
\vspace{0.5cm}

% 题目开始
\textbf{【单选】下列关于冯·诺依曼结构描述错误的是（\quad ）}
\begin{enumerate}[label=\Alph*., itemsep=2pt, parsep=0pt]
    \item 数据与程序采用十进制的方式进行存储 
    \item 最早出自《关于EDVAC的报告草案》 
    \item 采用“存储程序”的思想 
    \item 工作时自动从存储器中取出指令加以执行
\end{enumerate}
\vspace{0.4cm}

\textbf{【单选】下列关于冯·诺依曼结构计算机中英文对照关系错误的是（\quad ）}
\begin{enumerate}[label=\Alph*., itemsep=2pt, parsep=0pt]
    \item 运算器——CA 
    \item 控制器——CC 
    \item 管理器——M 
    \item 输入设备——I 
    \item 外部记录设备——R 
    \item 输出设备——O
\end{enumerate}
\vspace{0.4cm}

\textbf{【多选】下列哪些部件并称为计算机中的“大脑”？（\quad ）}
\begin{enumerate}[label=\Alph*., itemsep=2pt, parsep=0pt]
    \item 控制器 
    \item 运算器 
    \item 外部存储器 
    \item 主存储器 
    \item 寄存器
\end{enumerate}
\vspace{0.4cm}

\textbf{【多选】下列哪些设备属于输入设备？（\quad ）}
\begin{enumerate}[label=\Alph*., itemsep=2pt, parsep=0pt]
    \item 鼠标 
    \item 键盘 
    \item 显示器 
    \item 打印机
    \item 麦克风 
    \item 摄像头 
    \item 音箱
\end{enumerate}
\vspace{0.4cm}

\textbf{【多选】在传统台式机主板的南北桥结构中，下列设备中哪些属于南桥？（\quad ）}
\begin{enumerate}[label=\Alph*., itemsep=2pt, parsep=0pt]
    \item U盘
    \item 内存
    \item 键盘
    \item 鼠标
    \item PCIe显卡
    \item 硬盘
    \item CPU
\end{enumerate}
\vspace{0.4cm}


\textbf{【单选】在计算机结构的简化模型中，下列哪个寄存器是用于记录存储单元地址的？（\quad ）}
\begin{enumerate}[label=\Alph*., itemsep=2pt, parsep=0pt]
    \item IR 
    \item PC 
    \item MAR
    \item MDR
    \item R0
\end{enumerate}
\vspace{0.4cm}


\textbf{【多选】执行一条指令的全过程中，指令的编码会出现在CPU中的哪些部件？（\quad ）}
\begin{enumerate}[label=\Alph*., itemsep=2pt, parsep=0pt]
    \item R0
    \item IR
    \item MDR
    \item MAR
    \item PC
    \item ALU
\end{enumerate}
\vspace{0.6cm}

\newpage

\textbf{【问答题】课程中讲解的指令是“ADD R, M”，如果换成“ADD M, R”（即运算结果存入存储器），指令编码规则不变，即第0-1位表示操作码，第2-3位表示寄存器编号，第4-7位表示存储器单元地址。假设当前PC为0011，存储器状态如下：}

\begin{table}[htbp]
    \centering
    \begin{tabular}{cc}
        \toprule
        \textbf{地址} & \textbf{内容} \\
        \midrule
        0001 & 0100 0111 \\
        0010 & 0000 0101 \\
        0011 & 0000 0001 \\
        0100 & 0100 0010 \\
        0101 & 1100 0111 \\
        0110 & 0010 1010 \\
        0111 & 0001 0110 \\
        1000 & 1001 0001 \\
        \bottomrule
    \end{tabular}
\end{table}

\textbf{请回答以下问题：}
\begin{enumerate}[label=（\arabic*）, itemsep=4pt, parsep=0pt]
    \item 仅考虑已有信息，会运行几条指令？
    \item 每条指令的运行过程中，分别会读几次存储器、写几次存储器？
    \item 每条指令的运行结束后，哪些寄存器和存储单元发生了变化、发生了什么样的变化？
\end{enumerate}

\end{document}